% Test-case specification TC-0005 (CTest test TC-0005-filedata).
\documentclass[letterpaper]{article}

\usepackage[letterpaper,margin=1in]{geometry}
\emergencystretch=3em

% Compile from the repository root so this path resolves to
% the pinned WSP submodule.
\input{wsp/testing/test-case-library.tex}

\begin{document}

\def\TCID{TC-0005}
\def\TCTitle{%
  Print \#, Write \#, Input \#, Get \# and Put \#}
\def\TCPurpose{%
  Verify \texttt{Print \#}, \texttt{Write \#}, \texttt{Input \#},
  \texttt{Get \#} and \texttt{Put \#} against what a Visual Basic 6 program
  left in its files and variables.}
\def\TCPriority{\TCPriorityHigh}

\def\TCRequirementRef{%
  REQ-0003, REQ-0004}
\def\TCDesignRef{%
  Not applicable.}
\def\TCFeatureRef{%
  \texttt{tests/\allowbreak{}vbrun\_\allowbreak{}file\_\allowbreak{}data\_\allowbreak{}tests.\allowbreak{}cpp},
  \texttt{src/\allowbreak{}helpers\_\allowbreak{}openfiles.\allowbreak{}cpp},
  \texttt{tests/\allowbreak{}cases/\allowbreak{}fileprint.\allowbreak{}expected},
  \texttt{fileinput.\allowbreak{}expected},
  \texttt{filenumbers.\allowbreak{}expected}}
\def\TCTestTechnique{%
  Back-to-back comparison with recorded reference output.}

\def\TCPreconditions{%
  The library and the test program are built
  (\texttt{scripts/\allowbreak{}Invoke-Build.\allowbreak{}ps1}) and the test
  is registered with CTest.}
\def\TCEnvironment{%
  Windows with the toolchain of
  \texttt{scripts/\allowbreak{}Invoke-Build.\allowbreak{}ps1}; the local
  gate builds and runs x64 Debug, CI builds and runs the other
  architectures. The library is built from this repository in
  \texttt{out/\allowbreak{}build}.}
\def\TCAssumptions{%
  Not applicable.}

\def\TCInputData{%
  The statements of the experiments \texttt{filePrint} and
  \texttt{fileInput} of
  \texttt{scripts/\allowbreak{}vb6\_\allowbreak{}harness\_\allowbreak{}gen.\allowbreak{}py}
  and the number formats of
  \texttt{filenumbers\_\allowbreak{}values.\allowbreak{}inc}.}
\def\TCInitialState{%
  The working directory \texttt{work-filedata}, used by no other test.}

\def\TCProcedure{%
  \begin{enumerate}
    \item CTest runs
      \texttt{vbrun\_\allowbreak{}file\_\allowbreak{}data\_\allowbreak{}tests
      fileprint.\allowbreak{}expected fileinput.\allowbreak{}expected
      filenumbers.\allowbreak{}expected}.
    \item Run the print, input and number statement groups through the helpers.
    \item Record the file contents (as hexadecimal) and the variable values.
    \item Compare each group with its expected file.
    \item CTest records the exit code and, on failure, the program's output
      (\texttt{--output-on-failure}).
  \end{enumerate}}
\def\TCExpectedResult{%
  Each group's lines equal the expected file's lines; the program exits 0.}
\def\TCPostConditions{%
  No file or registry state remains that later tests depend on.}

\TCTemplate

\end{document}
